% Milestone 2 test documentation -- CRAN Polish Mirror Demand & GitHub Release Pulse Analysis System
% Build: python tests/e2e/make_tex.py && latexmk -pdf -outdir=build Milestone_2_Tests.tex   (output renamed TeamName_Tests.pdf)
\documentclass[11pt,a4paper]{article}

\usepackage[utf8]{inputenc}
\usepackage[T1]{fontenc}
\usepackage{lmodern}
\usepackage[margin=2.3cm]{geometry}
\usepackage{booktabs}
\usepackage{tabularx}
\usepackage{longtable}
\usepackage{array}
\usepackage{ragged2e}
\usepackage{enumitem}
\usepackage{xcolor}
\usepackage{float}
\usepackage[font=small,labelfont=bf]{caption}
\usepackage{listings}
\usepackage{fancyhdr}
\usepackage[colorlinks=true,linkcolor=blue!50!black,urlcolor=blue!50!black]{hyperref}

\setlist{itemsep=2pt,topsep=4pt}
\setlength{\parindent}{0pt}
\setlength{\parskip}{5pt}
\emergencystretch=3em
\renewcommand{\arraystretch}{1.15}
\setcounter{tocdepth}{1}
\newcolumntype{L}[1]{>{\RaggedRight\arraybackslash}p{#1}}
\newcolumntype{Y}{>{\RaggedRight\arraybackslash}X}
\captionsetup[table]{position=top,skip=4pt}

\pagestyle{fancy}
\fancyhf{}
\lhead{Big Data Analysis Classes | 2026Z}
\rhead{Milestone 2 -- Test Documentation}
\cfoot{\thepage}

\newcommand{\teamname}{TeamName}

\lstset{basicstyle=\ttfamily\footnotesize,breaklines=true,frame=single,columns=fullflexible,
  keepspaces=true,showstringspaces=false,inputencoding=utf8,extendedchars=true,
  literate={→}{{->}}1 {…}{{...}}1 {≥}{{>=}}1 {≤}{{<=}}1}

\title{\textbf{Project: CRAN Polish Mirror Demand \& GitHub Release Pulse Analysis System}\\[4pt]
\large Milestone 2 -- Test Documentation}
\author{Jakub Grzywaczewski, Mateusz Nizwantowski, Aleks Kapich, \L{}ukasz Grabarski}
\date{\today}

\begin{document}
\maketitle
\tableofcontents
\newpage

% =====================================================================
\section{Scope and approach}
% =====================================================================

This document covers the testing required for Milestone 2:

\begin{itemize}
  \item \textbf{Functional tests} of the \emph{data acquisition} modules (CRAN log source, GitHub Events poller, watch-list) and of the
        \emph{pre-processing} modules (the three Flink jobs that validate, deduplicate and transform the data);
  \item \textbf{Quality tests} of the \emph{data transformation} processes (conformance to schemas, completeness, uniqueness, partitioning, time bounds);
  \item \textbf{Tabular test case documentation} with test objective, test steps, expected result and actual result.
\end{itemize}

The tests are executable. A harness (\texttt{tests/e2e/run\_e2e.py} in the repository) runs against the real stack -- Kafka, Schema Registry,
HDFS and Flink -- and writes one log per test case. The \emph{actual result} columns and the printouts in the appendix are generated from that
run, and the log is the evidence, so screenshots are not needed. Analytical components that are still work in progress (Cassandra tables, Spark,
MLflow, Streamlit) are not tested here.

\subsection{What is tested, and how the data is chosen}

\begin{itemize}
  \item \textbf{Live data flows through the stack during the tests.} The log generator produces about 4 requests per second and the GitHub poller
        polls the real GitHub Events API with the project's token. Tests that check behaviour on this data (counts, schema, partitioning, API limits)
        therefore run on real events.
  \item \textbf{Injected test records.} For cases where the expected outcome must be known exactly (rejection reasons, deduplication, field mapping),
        the harness produces clearly marked records next to the live data: IP addresses from the documentation range \texttt{203.0.113.0/24},
        package names \texttt{e2etest.<tag>}, GitHub event ids \texttt{e2e-<tag>-*}. A random tag per run keeps runs apart.
  \item \textbf{Fresh stack.} Before the run the dev environment was wiped (\texttt{deploy/dev\_clean.sh}) and started from scratch
        (\texttt{deploy/up.sh}, \texttt{deploy/start\_logs.sh}, \texttt{deploy/start\_github.sh}), so the results do not depend on earlier data.
\end{itemize}

\textbf{Note on re-runs.} The full run (21 cases) was followed by a re-run of FP-02 and FP-05 with a new tag, because the first attempt failed on a defect
in the test harness (a DataFrame was used as a boolean), not in the pipeline; FA-03 and FA-04 were re-run to clean up the formatting of their summary line.
The logs and tables show the latest run of each case.

\subsection{Test environment}

\begin{table}[H]
\centering
\caption{Environment under test.}
\small
\begin{tabularx}{\textwidth}{@{}L{4.2cm}Y@{}}
\toprule
\textbf{Item} & \textbf{Value} \\
\midrule
Environment & dev, Docker Compose project \texttt{bigdata}, one host (8 CPU cores, 16\,GB RAM), Docker 29.7.2 / Compose 5.5.0 \\
Data platform & Kafka 4.3.2 (KRaft), Confluent Schema Registry 8.4.1, HDFS 3.4.2 (1 DataNode), Flink 1.20.3 (PyFlink), Cassandra 5.0.9, Airflow 3.4.0 \\
Data sources & Log generator (about 4 events/s, same format and topic as the future Filebeat input); GitHub Events API through the poller (watch-list: tier-1 organisations plus 100 other owners) \\
Test harness & Python 3.12, pytest, confluent-kafka, fastavro, pandas, pyarrow; HDFS read through \texttt{docker compose exec namenode hdfs dfs} \\
Test date & 23 November 2026 \\
\bottomrule
\end{tabularx}
\end{table}

\subsection{How to reproduce}

\begin{lstlisting}
deploy/dev_clean.sh -y && deploy/up.sh            # fresh dev stack
deploy/start_logs.sh && deploy/start_github.sh    # data sources (GITHUB_TOKEN in the environment)
pip install -r tests/e2e/requirements.txt
python tests/e2e/run_e2e.py                       # about 15 minutes; logs in deliverables/test_logs/
python tests/e2e/make_tex.py                      # regenerates the tables of this document
\end{lstlisting}

The run takes time because the Flink Parquet sinks finalise files on a 60\,s checkpoint, the poller needs several rounds after its restart, and
the completeness test waits until all consumers have caught up.

\subsection{Summary of results}

\begin{table}[H]
\centering
\caption{Results by group.}
\begin{tabular}{@{}lrrr@{}}
\toprule
\textbf{Group} & \textbf{Cases} & \textbf{Passed} & \textbf{Failed} \\
\midrule
Functional tests: data acquisition & 8 & 8 & 0 \\
Functional tests: pre-processing & 6 & 6 & 0 \\
Quality tests: data transformation & 7 & 7 & 0 \\
\midrule
\textbf{Total} & 21 & 21 & 0 \\
\bottomrule
\end{tabular}

\end{table}

\subsection{Traceability}

\begin{table}[H]
\centering
\caption{Modules under test and the cases that cover them.}
\small
\begin{tabularx}{\textwidth}{@{}L{4.6cm}L{2.6cm}Y@{}}
\toprule
\textbf{Module} & \textbf{Cases} & \textbf{What is checked} \\
\midrule
CRAN log source (generator; Filebeat later) & FA-01 & Delivery of well-formed common-log lines to Kafka \\
Platform initialisation (\texttt{deploy/init.sh}) & FA-02 & Topics, Avro subjects, HDFS directories \\
GitHub poller (code) & FA-03, FA-04 & Unit tests with mocked HTTP; live tests against the real API \\
Watch-list build & FA-05 & Watch-list from crantogit, publication to \texttt{github\_repo\_list} \\
GitHub poller (runtime) & FA-06, FA-07, FA-08 & Message format, ETag and quota handling, restart safety \\
Flink job \texttt{cran\_logs\_ingest} & FP-01, FP-02, FP-06 & Parsing, validation and rejection, raw storage \\
Flink job \texttt{cran\_package\_downloads} & FP-03 & Extraction of package downloads \\
Flink job \texttt{github\_events\_ingest} & FP-04, FP-05, FP-06 & Mapping, deduplication, validation, raw storage \\
All transformations & Q-01 to Q-07 & Schemas, uniqueness, completeness, partitioning, keys, rejections, time bounds \\
\bottomrule
\end{tabularx}
\end{table}

% =====================================================================
% \newgeometry cannot change the paper size, so the text block is given explicitly and the PDF page size is set directly.
\newgeometry{left=2cm,top=2cm,textwidth=25.7cm,textheight=16.2cm,headheight=14pt,headsep=0.5cm,footskip=0.9cm}
\pdfpagewidth=29.7cm \pdfpageheight=21cm
\section{Functional tests: data acquisition}
% =====================================================================
Test case IDs \texttt{FA-nn}. Status and actual results come from the run of 23 November 2026; the full printout of each case is in the appendix.

{\footnotesize
\begin{longtable}{@{}L{1.1cm}L{5.8cm}L{5.4cm}L{3.9cm}L{5.5cm}L{1.3cm}@{}}
\toprule
\textbf{ID} & \textbf{Test objective} & \textbf{Test steps} & \textbf{Expected result} & \textbf{Actual result} & \textbf{Status} \\
\midrule
\endfirsthead
\toprule
\textbf{ID} & \textbf{Test objective} & \textbf{Test steps} & \textbf{Expected result} & \textbf{Actual result} & \textbf{Status} \\
\midrule
\endhead
\bottomrule
\endfoot
\textbf{FA-01} & \textbf{Log source (log generator; Filebeat later) delivers log lines to Kafka}\newline Check that the current CRAN log source, the log generator that simulates the mirror until Filebeat is deployed, delivers well-formed access-log lines to the raw topic. It tests the interface (topic, line format) that Filebeat must also meet, not Filebeat itself. & 1.~Read the end offsets of cran\_\allowbreak{}mirror\_\allowbreak{}logs, wait 10 s, read them again.\newline 2.~Read a sample of messages and match each against the Apache common log format. & At least 20 new messages in 10 s; every sampled message is a valid common-log line. & 39 messages in 10 s; 200/\allowbreak{}200 sampled lines are valid common-log lines.\newline \textit{Log: Appendix, FA-01.} & \textbf{\textcolor{green!45!black}{PASS}} \\
\midrule[0.3pt]
\textbf{FA-02} & \textbf{Platform initialisation: topics, schemas, HDFS directories}\newline Check that deploy/\allowbreak{}init.sh created everything the acquisition modules write to. & 1.~Describe the Kafka topics.\newline 2.~List the Schema Registry subjects.\newline 3.~List the HDFS raw and clean directories. & Six topics with 3 partitions, three Avro subjects and eight HDFS directories exist. & 6 topics (3 partitions), 3 Avro subjects and 8 HDFS data directories exist.\newline \textit{Log: Appendix, FA-02.} & \textbf{\textcolor{green!45!black}{PASS}} \\
\midrule[0.3pt]
\textbf{FA-03} & \textbf{GitHub poller unit tests (mocked HTTP)}\newline Check the poller logic offline: ETags, rate limits, scheduling, retries, watch-list, Kafka sink, state. & 1.~Run pytest -m 'not live' in poller/\allowbreak{}. & All tests pass. & pytest -m 'not live': 107 passed, 5 deselected in 3.56s.\newline \textit{Log: Appendix, FA-03.} & \textbf{\textcolor{green!45!black}{PASS}} \\
\midrule[0.3pt]
\textbf{FA-04} & \textbf{GitHub poller live tests (real API)}\newline Check the poller against the real GitHub Events API with the project token. & 1.~Run pytest -m live in poller/\allowbreak{} with GITHUB\_\allowbreak{}TOKEN set. & All live tests pass. & pytest -m live (real GitHub API): 5 passed, 107 deselected in 7.74s.\newline \textit{Log: Appendix, FA-04.} & \textbf{\textcolor{green!45!black}{PASS}} \\
\midrule[0.3pt]
\textbf{FA-05} & \textbf{Watch-list build from crantogit and publication to Kafka}\newline Check that the watch-list is built from crantogit.csv and published to github\_\allowbreak{}repo\_\allowbreak{}list. & 1.~Read /\allowbreak{}data/\allowbreak{}watchlist.csv from the poller container as a DataFrame.\newline 2.~Validate columns, tiers and endpoints.\newline 3.~Read github\_\allowbreak{}repo\_\allowbreak{}list and compare the rows of the latest build. & Valid watch-list; the latest build published exactly the rows of the file. & 2627 watch-list rows (\{2: 1983, 1: 644\}), all published to github\_\allowbreak{}repo\_\allowbreak{}list.\newline \textit{Log: Appendix, FA-05.} & \textbf{\textcolor{green!45!black}{PASS}} \\
\midrule[0.3pt]
\textbf{FA-06} & \textbf{GitHub events reach Kafka with the documented format}\newline Check the structure, keys and package matching of the events written by the poller. & 1.~Read all messages of github\_\allowbreak{}repo\_\allowbreak{}updates.\newline 2.~Check the keys, mandatory fields, track values and package matching. & Key = repository name; every message has the documented fields; events are matched to CRAN packages. & 4814 events in Kafka, keys correct, 3811 matched to CRAN packages.\newline \textit{Log: Appendix, FA-06.} & \textbf{\textcolor{green!45!black}{PASS}} \\
\midrule[0.3pt]
\textbf{FA-07} & \textbf{API limit handling: conditional requests and quota}\newline Check from the poller's own request log that ETags are used and the quota is respected. & 1.~Read the request log of the poller (requests-*.jsonl) as a DataFrame.\newline 2.~Count polls by HTTP status.\newline 3.~Check rate-limit flags and the remaining quota. & Both 200 and 304 responses occur; no rate limit; remaining quota stays above the reserve. & 588 polls, statuses \{200: 554, 304: 34\}; no rate limit; minimum remaining quota 4303.\newline \textit{Log: Appendix, FA-07.} & \textbf{\textcolor{green!45!black}{PASS}} \\
\midrule[0.3pt]
\textbf{FA-08} & \textbf{Poller restart does not duplicate events}\newline Check that the poller keeps its state (ETags, written event ids) across a restart. & 1.~Record the event ids in github\_\allowbreak{}repo\_\allowbreak{}updates.\newline 2.~Restart the poller container and wait for several poll rounds.\newline 3.~Check in the poller's request log that it polled again and that polls after the restart got 304 (ETags kept).\newline 4.~Check that no event id occurs twice and earlier events are still present. & Polls after the restart; some answered 304; no duplicate event ids. & 214 polls after the restart (\{200: 205, 304: 9\}); 8031 distinct events, 0 duplicates.\newline \textit{Log: Appendix, FA-08.} & \textbf{\textcolor{green!45!black}{PASS}} \\
\midrule[0.3pt]
\end{longtable}}


% =====================================================================
\section{Functional tests: pre-processing}
% =====================================================================
Test case IDs \texttt{FP-nn}. Records marked \texttt{e2etest.<tag>} / \texttt{203.0.113.x} are injected by the harness (see the scope section).

{\footnotesize
\begin{longtable}{@{}L{1.1cm}L{5.8cm}L{5.4cm}L{3.9cm}L{5.5cm}L{1.3cm}@{}}
\toprule
\textbf{ID} & \textbf{Test objective} & \textbf{Test steps} & \textbf{Expected result} & \textbf{Actual result} & \textbf{Status} \\
\midrule
\endfirsthead
\toprule
\textbf{ID} & \textbf{Test objective} & \textbf{Test steps} & \textbf{Expected result} & \textbf{Actual result} & \textbf{Status} \\
\midrule
\endhead
\bottomrule
\endfoot
\textbf{FP-01} & \textbf{Valid CRAN log lines are parsed into the clean topic}\newline Check the parsing of the Flink log job for valid lines. & 1.~Produce six valid test lines to cran\_\allowbreak{}mirror\_\allowbreak{}logs.\newline 2.~Read them from cran\_\allowbreak{}mirror\_\allowbreak{}logs\_\allowbreak{}clean (Avro).\newline 3.~Compare every field with the input. & All six lines arrive with correct fields; '-' size becomes null. & All 6 valid test lines reached cran\_\allowbreak{}mirror\_\allowbreak{}logs\_\allowbreak{}clean with correct fields; '-' size became null.\newline \textit{Log: Appendix, FP-01.} & \textbf{\textcolor{green!45!black}{PASS}} \\
\midrule[0.3pt]
\textbf{FP-02} & \textbf{Invalid CRAN log lines are rejected with a reason}\newline Check that invalid lines are kept in the bad zone and never reach the clean data. & 1.~Produce six invalid lines (garbage, status 999, method FOO, +1 day, -30 days, offset +0200).\newline 2.~Read /\allowbreak{}data/\allowbreak{}raw/\allowbreak{}cran\_\allowbreak{}mirror\_\allowbreak{}logs\_\allowbreak{}bad as a DataFrame.\newline 3.~Compare reject\_\allowbreak{}reason with the expectation.\newline 4.~Check the clean topic. & Each line is in the bad zone with the expected reason and absent from the clean topic. & 6/\allowbreak{}6 invalid lines were rejected with the expected reject\_\allowbreak{}reason and none reached the clean topic.\newline \textit{Log: Appendix, FP-02.} & \textbf{\textcolor{green!45!black}{PASS}} \\
\midrule[0.3pt]
\textbf{FP-03} & \textbf{Package downloads are extracted from the clean log}\newline Check the download extraction job: package, version, artifact type, target system, outcome. & 1.~Use the valid test lines of FP-01 (source, macOS, Windows, a 404, PACKAGES.rds, an HTML page).\newline 2.~Read cran\_\allowbreak{}mirror\_\allowbreak{}package\_\allowbreak{}downloads (Avro).\newline 3.~Compare the derived fields; check the skipped requests. & Four downloads with correct fields (404 = FAILURE); the two non-package requests are skipped. & 4/\allowbreak{}4 package requests were extracted correctly (source, macos, windows, 404 as FAILURE); 2 non-downloads skipped.\newline \textit{Log: Appendix, FP-03.} & \textbf{\textcolor{green!45!black}{PASS}} \\
\midrule[0.3pt]
\textbf{FP-04} & \textbf{GitHub events: mapping and deduplication}\newline Check the GitHub Flink job: field extraction, nullable fields and removal of repeated event ids. & 1.~Produce a ReleaseEvent twice and a PushEvent without organisation to github\_\allowbreak{}repo\_\allowbreak{}updates.\newline 2.~Read github\_\allowbreak{}repo\_\allowbreak{}updates\_\allowbreak{}clean (Avro).\newline 3.~Compare the fields; count the occurrences of the repeated event. & The repeated event appears once; fields are mapped; missing organisation is null. & Duplicate event kept once; release\_\allowbreak{}tag, ref, packages and nullable org\_\allowbreak{}login mapped correctly.\newline \textit{Log: Appendix, FP-04.} & \textbf{\textcolor{green!45!black}{PASS}} \\
\midrule[0.3pt]
\textbf{FP-05} & \textbf{Invalid GitHub records are rejected with a reason}\newline Check the validation rules of the GitHub job. & 1.~Produce five invalid records (30 days old, track 'other', not JSON, no repo name, invalid timestamp).\newline 2.~Read /\allowbreak{}data/\allowbreak{}raw/\allowbreak{}github\_\allowbreak{}repo\_\allowbreak{}updates\_\allowbreak{}bad as a DataFrame.\newline 3.~Compare reject\_\allowbreak{}reason; check the clean topic. & Each record is in the bad zone with the expected reason and absent from the clean topic. & 5/\allowbreak{}5 invalid GitHub events were rejected with the expected reject\_\allowbreak{}reason.\newline \textit{Log: Appendix, FP-05.} & \textbf{\textcolor{green!45!black}{PASS}} \\
\midrule[0.3pt]
\textbf{FP-06} & \textbf{Raw records are stored unchanged}\newline Check that the raw zone keeps every input record byte for byte, valid or not. & 1.~Read the text files of /\allowbreak{}data/\allowbreak{}raw/\allowbreak{}cran\_\allowbreak{}mirror\_\allowbreak{}logs and /\allowbreak{}data/\allowbreak{}raw/\allowbreak{}github\_\allowbreak{}repo\_\allowbreak{}updates.\newline 2.~Look for each injected record. & All injected records are present unchanged. & All 12 injected log lines and 8 injected GitHub records are stored unchanged in the raw zone.\newline \textit{Log: Appendix, FP-06.} & \textbf{\textcolor{green!45!black}{PASS}} \\
\midrule[0.3pt]
\end{longtable}}


% =====================================================================
\section{Quality tests: data transformation}
% =====================================================================
Test case IDs \texttt{Q-nn}. These cases run after the producers were stopped and the Flink consumers had caught up, so that counts are stable.

{\footnotesize
\begin{longtable}{@{}L{1.1cm}L{5.8cm}L{5.4cm}L{3.9cm}L{5.5cm}L{1.3cm}@{}}
\toprule
\textbf{ID} & \textbf{Test objective} & \textbf{Test steps} & \textbf{Expected result} & \textbf{Actual result} & \textbf{Status} \\
\midrule
\endfirsthead
\toprule
\textbf{ID} & \textbf{Test objective} & \textbf{Test steps} & \textbf{Expected result} & \textbf{Actual result} & \textbf{Status} \\
\midrule
\endhead
\bottomrule
\endfoot
\textbf{Q-01} & \textbf{Schema conformance and value constraints}\newline Check that the transformed data conform to the Avro schemas and the constraints Avro cannot express. & 1.~Decode every message of the three derived topics with the schema from the registry.\newline 2.~Compare the Parquet columns with the Avro fields.\newline 3.~Check status range, methods, allowed values. & Everything decodes; columns match; all constraints hold. & All three clean topics decode with the registry schema; Parquet columns match; value constraints hold.\newline \textit{Log: Appendix, Q-01.} & \textbf{\textcolor{green!45!black}{PASS}} \\
\midrule[0.3pt]
\textbf{Q-02} & \textbf{No duplicate GitHub events in the clean data}\newline Check that deduplication on event\_\allowbreak{}id works on the live data. & 1.~Count distinct event ids in the clean topic and in the Parquet files.\newline 2.~Compare with the raw topic. & No event id occurs twice in the clean data. & 2905 clean events, 0 duplicates (topic and Parquet); 1 repeated raw records were dropped.\newline \textit{Log: Appendix, Q-02.} & \textbf{\textcolor{green!45!black}{PASS}} \\
\midrule[0.3pt]
\textbf{Q-03} & \textbf{Completeness: no record is lost between raw, clean and rejected}\newline Check that raw = clean + rejected for the CRAN logs, and that Kafka and HDFS agree. & 1.~Stop the producers and wait until the Flink consumer lag is 0 and the files are finalised.\newline 2.~Count the raw topic, raw text, clean topic, clean Parquet and bad Parquet. & raw = clean + rejected; topic counts equal file counts. & 2462 raw lines = 2456 clean + 6 rejected; Kafka and HDFS counts agree.\newline \textit{Log: Appendix, Q-03.} & \textbf{\textcolor{green!45!black}{PASS}} \\
\midrule[0.3pt]
\textbf{Q-04} & \textbf{Partitioning by event time}\newline Check that Parquet files are in the dt/\allowbreak{}hr directory of the record's time. & 1.~Read the Parquet datasets and compare dt and hr with event\_\allowbreak{}ts (arrival\_\allowbreak{}ts for rejects). & Every row is in the correct partition. & All Parquet rows are in the partition matching their event time (rejected lines: arrival time).\newline \textit{Log: Appendix, Q-04.} & \textbf{\textcolor{green!45!black}{PASS}} \\
\midrule[0.3pt]
\textbf{Q-05} & \textbf{Keys and partition locality}\newline Check that records of one package or repository stay in one partition. & 1.~Group the downloads topic by package\_\allowbreak{}name and the GitHub topic by repo\_\allowbreak{}name.\newline 2.~Count partitions per key. & Exactly one partition per key; key equals the field. & Records with the same package\_\allowbreak{}name /\allowbreak{} repo\_\allowbreak{}name always go to the same partition, so their order is kept.\newline \textit{Log: Appendix, Q-05.} & \textbf{\textcolor{green!45!black}{PASS}} \\
\midrule[0.3pt]
\textbf{Q-06} & \textbf{Reject accounting}\newline Check that rejections happen only for the documented reasons and only where expected. & 1.~Count the rejected CRAN lines by reason and compare them with the injected invalid lines.\newline 2.~Count the rejected GitHub records by reason. & CRAN rejects are exactly the injected ones; GitHub reasons are documented, stale history is rejected. & CRAN: 6 rejects, all of them injected test lines. GitHub: 2219 rejects (\{'stale\_\allowbreak{}timestamp': 2215, 'invalid\_\allowbreak{}track': 1, 'unparseable': 1, 'missing\_\allowbreak{}field': 1, 'invalid\_\allowbreak{}timestamp': 1\}), 2905 clean events.\newline \textit{Log: Appendix, Q-06.} & \textbf{\textcolor{green!45!black}{PASS}} \\
\midrule[0.3pt]
\textbf{Q-07} & \textbf{Event-time plausibility of the clean data}\newline Check the time bounds the ingest jobs promise. & 1.~Compute event\_\allowbreak{}ts - arrival\_\allowbreak{}ts (CRAN) and fetched\_\allowbreak{}ts - event\_\allowbreak{}ts (GitHub) over all clean rows. & CRAN: -7 d to +5 min; GitHub: 0 to about 7 d. & All clean records satisfy the event-time bounds enforced by the ingest jobs.\newline \textit{Log: Appendix, Q-07.} & \textbf{\textcolor{green!45!black}{PASS}} \\
\midrule[0.3pt]
\end{longtable}}

\clearpage
\restoregeometry
\pdfpagewidth=21cm \pdfpageheight=29.7cm

% =====================================================================
\section{Example data}
% =====================================================================

The following printout shows the transformed data as read from the Parquet files in HDFS into pandas DataFrames: shape, column types and the
first rows of each dataset, plus a few aggregates. It is output of the same run as the test cases.

\lstinputlisting[basicstyle=\ttfamily\tiny,breaklines=true]{test_logs/DATA-EXAMPLES.log}

% =====================================================================
\appendix
\section{Test logs}
% =====================================================================

Complete printout of every test case: objective, steps, expected result, observed output and result.

\subsection*{FA-01 -- Log source (log generator; Filebeat later) delivers log lines to Kafka}
\lstinputlisting[basicstyle=\ttfamily\tiny,breaklines=true,breakatwhitespace=false]{test_logs/FA-01.log}

\subsection*{FA-02 -- Platform initialisation: topics, schemas, HDFS directories}
\lstinputlisting[basicstyle=\ttfamily\tiny,breaklines=true,breakatwhitespace=false]{test_logs/FA-02.log}

\subsection*{FA-03 -- GitHub poller unit tests (mocked HTTP)}
\lstinputlisting[basicstyle=\ttfamily\tiny,breaklines=true,breakatwhitespace=false]{test_logs/FA-03.log}

\subsection*{FA-04 -- GitHub poller live tests (real API)}
\lstinputlisting[basicstyle=\ttfamily\tiny,breaklines=true,breakatwhitespace=false]{test_logs/FA-04.log}

\subsection*{FA-05 -- Watch-list build from crantogit and publication to Kafka}
\lstinputlisting[basicstyle=\ttfamily\tiny,breaklines=true,breakatwhitespace=false]{test_logs/FA-05.log}

\subsection*{FA-06 -- GitHub events reach Kafka with the documented format}
\lstinputlisting[basicstyle=\ttfamily\tiny,breaklines=true,breakatwhitespace=false]{test_logs/FA-06.log}

\subsection*{FA-07 -- API limit handling: conditional requests and quota}
\lstinputlisting[basicstyle=\ttfamily\tiny,breaklines=true,breakatwhitespace=false]{test_logs/FA-07.log}

\subsection*{FA-08 -- Poller restart does not duplicate events}
\lstinputlisting[basicstyle=\ttfamily\tiny,breaklines=true,breakatwhitespace=false]{test_logs/FA-08.log}

\subsection*{FP-01 -- Valid CRAN log lines are parsed into the clean topic}
\lstinputlisting[basicstyle=\ttfamily\tiny,breaklines=true,breakatwhitespace=false]{test_logs/FP-01.log}

\subsection*{FP-02 -- Invalid CRAN log lines are rejected with a reason}
\lstinputlisting[basicstyle=\ttfamily\tiny,breaklines=true,breakatwhitespace=false]{test_logs/FP-02.log}

\subsection*{FP-03 -- Package downloads are extracted from the clean log}
\lstinputlisting[basicstyle=\ttfamily\tiny,breaklines=true,breakatwhitespace=false]{test_logs/FP-03.log}

\subsection*{FP-04 -- GitHub events: mapping and deduplication}
\lstinputlisting[basicstyle=\ttfamily\tiny,breaklines=true,breakatwhitespace=false]{test_logs/FP-04.log}

\subsection*{FP-05 -- Invalid GitHub records are rejected with a reason}
\lstinputlisting[basicstyle=\ttfamily\tiny,breaklines=true,breakatwhitespace=false]{test_logs/FP-05.log}

\subsection*{FP-06 -- Raw records are stored unchanged}
\lstinputlisting[basicstyle=\ttfamily\tiny,breaklines=true,breakatwhitespace=false]{test_logs/FP-06.log}

\subsection*{Q-01 -- Schema conformance and value constraints}
\lstinputlisting[basicstyle=\ttfamily\tiny,breaklines=true,breakatwhitespace=false]{test_logs/Q-01.log}

\subsection*{Q-02 -- No duplicate GitHub events in the clean data}
\lstinputlisting[basicstyle=\ttfamily\tiny,breaklines=true,breakatwhitespace=false]{test_logs/Q-02.log}

\subsection*{Q-03 -- Completeness: no record is lost between raw, clean and rejected}
\lstinputlisting[basicstyle=\ttfamily\tiny,breaklines=true,breakatwhitespace=false]{test_logs/Q-03.log}

\subsection*{Q-04 -- Partitioning by event time}
\lstinputlisting[basicstyle=\ttfamily\tiny,breaklines=true,breakatwhitespace=false]{test_logs/Q-04.log}

\subsection*{Q-05 -- Keys and partition locality}
\lstinputlisting[basicstyle=\ttfamily\tiny,breaklines=true,breakatwhitespace=false]{test_logs/Q-05.log}

\subsection*{Q-06 -- Reject accounting}
\lstinputlisting[basicstyle=\ttfamily\tiny,breaklines=true,breakatwhitespace=false]{test_logs/Q-06.log}

\subsection*{Q-07 -- Event-time plausibility of the clean data}
\lstinputlisting[basicstyle=\ttfamily\tiny,breaklines=true,breakatwhitespace=false]{test_logs/Q-07.log}


\end{document}
